\documentclass[12pt]{article}
\usepackage[margin=1in]{geometry}
\usepackage{amsmath,amssymb,amsthm,mathrsfs}
\usepackage{hyperref}
\usepackage{enumitem}
\usepackage{booktabs}
\usepackage{caption}

\theoremstyle{plain}
\newtheorem{theorem}{Theorem}[section]
\newtheorem{proposition}[theorem]{Proposition}
\newtheorem{lemma}[theorem]{Lemma}
\newtheorem{corollary}[theorem]{Corollary}

\theoremstyle{definition}
\newtheorem{definition}[theorem]{Definition}
\newtheorem{remark}[theorem]{Remark}

\numberwithin{equation}{section}

\newcommand{\R}{\mathbb{R}}
\newcommand{\C}{\mathbb{C}}
\newcommand{\Z}{\mathbb{Z}}
\newcommand{\br}{\mathrm{br}}
\newcommand{\vis}{\mathrm{vis}}
\newcommand{\gen}{\mathrm{gen}}
\newcommand{\kker}{\mathrm{ker}}
\newcommand{\eff}{\mathrm{eff}}
\newcommand{\obs}{\mathrm{obs}}
\newcommand{\GUT}{\mathrm{GUT}}
\newcommand{\SM}{\mathrm{SM}}
\DeclareMathOperator{\sech}{sech}
\newcommand{\fU}{\mathfrak{U}}
\newcommand{\fI}{\mathfrak{I}}
\newcommand{\fV}{\mathfrak{V}}
\newcommand{\fB}{\mathfrak{B}}
\newcommand{\red}{\mathrm{red}}

\begin{document}

%% ============================================================
%%  TITLE
%% ============================================================
\begin{titlepage}
\centering
\vspace*{1cm}

{\LARGE\bfseries Standard Model Structure, GUT Scale, and Dark Sector\\[6pt]
from Elliptic Modular Geometry:\\[6pt]
Predictions of LoNalogy}

\vspace{1.5cm}

{\large Masayoshi Nakamura$^{1,*}$}

\vspace{0.8cm}

{\normalsize
$^1$\,Kitayono Mental Clinic,
338-0002, 4-20-14 Shimoochiai, Chuo-ku,\\
Saitama City, Saitama, Japan.
\texttt{mjolnir501@gmail.com}\\[4pt]
$^*$\,Corresponding author.
}

\vspace{0.8cm}

{\small
\noindent\textbf{Related deposits:}\\
LoNalogy framework (v87.1): \href{https://doi.org/10.5281/zenodo.19115338}{doi:10.5281/zenodo.19115338}\\
Yang--Mills proof (v90): \href{https://doi.org/10.5281/zenodo.19183838}{doi:10.5281/zenodo.19183838}
}

\vspace{1.2cm}

\begin{abstract}
\noindent
We derive Standard Model structure, GUT-scale predictions, and dark sector
properties from the LoNalogy framework---a constructive approach based on
the para-complex unit $\jmath^2=+1$ and the level-2 modular curve
$X(2)=\Gamma(2)\backslash\mathbb{H}$.
The gauge group $\mathrm{SU}(3)\times\mathrm{SU}(2)\times\mathrm{U}(1)$
emerges from the $3+2$ split of a rank-5 fiber bundle, forced by Yukawa
closure and anomaly cancellation.  Exactly three generations arise from
the three cusps $\{0,1,\infty\}$ of $X(2)$ via a cusp saturation
principle, with UV family symmetry $S_3$.  The GUT-scale gauge boson mass
is derived as
\[
M_X = \Lambda_*\!\left(\frac{r}{s}\right)^{1/3}Q_{\ker}^{-14/3}
= 3.455\times10^{15}\,\mathrm{GeV},
\quad (r,s)=(3,2),
\]
giving proton lifetime $\tau_p=1.49\times10^{38}\,\mathrm{yr}$, safely
beyond current Super-Kamiokande bounds.
The cosmological constant is reproduced as
$\Lambda_\eff=\Lambda_*^4 Q_{\ker}^{20}=2.94\times10^{-47}\,\mathrm{GeV}^4$,
within 17\% of the PDG 2025 value without free parameters.
The dark ($e_-$) sector, corresponding to $s_{\rm int}=4$ in the integer
chain, predicts a hidden fermion mass window $m_\chi\simeq338$--$348\,\mathrm{GeV}$
with direct-detection cross section $\sigma_p^{\rm SI}\sim10^{-48}\,\mathrm{cm}^2$,
at the frontier of current LZ/XENONnT sensitivity.
Halo density profiles are analytically core-like (log-slope $\approx-0.17$),
alleviating the core-cusp tension.
\end{abstract}
\end{titlepage}

\setcounter{page}{1}
\tableofcontents
\newpage

%% ============================================================
\section*{Introduction}
\addcontentsline{toc}{section}{Introduction}
%% ============================================================

The LoNalogy framework derives from two algebraic primitives: the standard
complex unit $i^2=-1$ and the para-complex unit $\jmath^2=+1$, $\jmath\neq\pm1$.
These define two idempotents
\[
e_+ = \frac{1+\jmath}{2},\qquad e_- = \frac{1-\jmath}{2},\qquad e_+e_-=0,
\]
which split any LoNalogy field into a \emph{visible} ($e_+$) and a
\emph{dark} ($e_-$) sector.  The modular geometry of $X(2)$ governs both
sectors simultaneously.

The companion paper \cite{YMmassGap} establishes the Yang--Mills mass gap
for every compact simple gauge group $G$ using this framework.  The present
paper extracts the particle-physics and cosmological predictions that emerge
from the same geometric data.

All numerical values in this paper are derived from a single datum:
$\nu_\br=12$ (the broken orbit dimension for the $(3,2)$-split of $\mathrm{SU}(5)$).
No additional free parameters are introduced.

%% ============================================================
\part{Standard Model Structure}
%% ============================================================

\section{Gauge Group from the $3+2$ Split}\label{sec:gaugegroup}

\subsection{Fiber bundle setup}

\begin{definition}[Rank-5 visible fiber]
The visible fiber bundle has structure group $\mathrm{SU}(5)$ with fiber
\[
V_{\rm fib} = L_1\oplus L_2\oplus L_3\oplus M_1\oplus M_2
=: V_3\oplus V_2,\qquad \det V_{\rm fib}\simeq\mathcal{O}.
\]
The $3+2$ split is forced by Yukawa closure (Theorem~\ref{thm:yukawa})
and anomaly cancellation (Proposition~\ref{prop:anomaly}).
\end{definition}

\begin{theorem}[Gauge algebra decomposition]\label{thm:gaugedecomp}
The $3+2$ holonomy split gives
\[
\mathfrak{su}(5)\to\mathfrak{su}(3)\oplus\mathfrak{su}(2)\oplus\mathfrak{u}(1).
\]
The off-diagonal bridge is
\[
W_X = \operatorname{Hom}(V_2,V_3)\oplus\overline{\operatorname{Hom}(V_2,V_3)}
\cong (\mathbf{3},\mathbf{2})_{-5/6}\oplus(\bar{\mathbf{3}},\mathbf{2})_{+5/6}.
\]
The low-energy visible algebra is precisely that of the Standard Model.
\end{theorem}

\begin{proof}
Decompose the $\mathfrak{su}(5)$ adjoint $\mathbf{24}$ under the block
structure $V_3\oplus V_2$:
\[
\mathbf{24}\to(\mathbf{8},\mathbf{1})_0\oplus(\mathbf{1},\mathbf{3})_0
\oplus(\mathbf{1},\mathbf{1})_0\oplus(\mathbf{3},\mathbf{2})_{-5/6}
\oplus(\bar{\mathbf{3}},\mathbf{2})_{+5/6}.
\]
The diagonal blocks give $\mathfrak{su}(3)\oplus\mathfrak{su}(2)\oplus\mathfrak{u}(1)$.
The off-diagonal $W_X$ are heavy ($M_X\sim10^{15}\,\mathrm{GeV}$)
and decouple at low energies.
\end{proof}

\begin{remark}
This is not an assumption of $\mathrm{SU}(5)$ with SM embedded by hand.
The rank-5 selection is forced by two independent constraints
(Yukawa cubic closure and gauge anomaly cancellation), and the $3+2$
holonomy split then uniquely determines the low-energy gauge algebra.
\end{remark}

\subsection{Rank-5 selection}\label{sec:rank5}

\begin{theorem}[Yukawa cubic closure]\label{thm:yukawa}
The natural map $\Lambda^2 V\otimes\Lambda^2 V\otimes V\to\det V$
is a determinant-valued $\mathrm{SU}(N)$-invariant if and only if $N=5$.
\end{theorem}

\begin{proof}
For $V=\mathbb{C}^N$, the exterior product defines a multilinear map
$\phi:\Lambda^2 V\otimes\Lambda^2 V\otimes V\to\Lambda^5 V$ by
$(u_1\wedge u_2)\otimes(v_1\wedge v_2)\otimes w\mapsto u_1\wedge u_2\wedge v_1\wedge v_2\wedge w$.
This map is $\mathrm{SU}(N)$-equivariant.  The target $\Lambda^5 V$ is one-dimensional
(isomorphic to $\det V$) if and only if $\dim V=5$.

For $N<5$: $\Lambda^5 V=0$, so $\phi=0$ trivially.
For $N=5$: $\Lambda^5 V\cong\mathbb{C}$ and $\phi\neq0$ since $\det g=1$ for $g\in\mathrm{SU}(5)$.
For $N>5$: $\Lambda^5 V$ has dimension $\binom{N}{5}>1$, so $\phi$ does not map into $\det V$.
\end{proof}

\begin{proposition}[Anomaly cancellation]\label{prop:anomaly}
For $\mathrm{SU}(N)$ with fermion representations $\Lambda^2 V$ and $V^*$:
\[
A(\Lambda^2 V)+A(V^*) = (N-4)+(-1) = N-5.
\]
Anomaly cancellation requires $N=5$.
\end{proposition}

\begin{proof}
The anomaly coefficient of a representation $R$ is defined by
$\mathrm{Tr}_R(T^a\{T^b,T^c\})=A(R)\,d^{abc}$ where $d^{abc}$ is the symmetric $d$-symbol.
By standard representation theory: $A(V)=1$ (fundamental, by normalization);
$A(V^*)=-A(V)=-1$ (conjugate); $A(\Lambda^2 V)=N-4$ (antisymmetric tensor, by the
index formula).  Thus $A(\Lambda^2 V)+A(V^*)=(N-4)-1=N-5$.
For quantum consistency the total anomaly must vanish: $N=5$.
\end{proof}

\begin{corollary}
The conjunction of Yukawa closure and anomaly cancellation uniquely forces
$N=5$.  Neither condition alone suffices; both independently select $N=5$.
\end{corollary}

\begin{proof}
Theorem~\ref{thm:yukawa} shows the Yukawa map is determinant-valued iff $N=5$.
Proposition~\ref{prop:anomaly} shows anomaly cancellation forces $N=5$ independently.
The conjunction thus uniquely selects $N=5$; no other value satisfies both constraints simultaneously.
\end{proof}

%% ============================================================
\section{Three Generations from Three Cusps}\label{sec:generations}
%% ============================================================

\begin{proposition}[Cusp count of $X(2)$]
The compactification $\bar{X}(2)\simeq\mathbb{P}^1$ has exactly three cusps:
\[
\{0,\ 1,\ \infty\}\subset\mathbb{P}^1.
\]
These are the degeneration loci of the Legendre family $E_\lambda$.
\end{proposition}

\begin{proof}
The modular curve $X(2)=\Gamma(2)\backslash\overline{\mathbb{H}}$ has genus 0 by the
Riemann--Hurwitz formula: $\Gamma(2)$ has index 6 in $\mathrm{SL}_2(\mathbb{Z})$,
giving $2g-2=6(2\cdot0-2)+(\text{ramification})$.  The cusps of $\Gamma(2)$ are the
orbits of $\mathbb{P}^1(\mathbb{Q})$ under $\Gamma(2)$: these are precisely the three
orbits $\{0\}$, $\{1\}$, $\{\infty\}$, corresponding to the three degeneration limits
$\lambda\to0$, $\lambda\to1$, $\lambda\to\infty$ of the Legendre family $y^2=x(x-1)(x-\lambda)$.
\end{proof}

\begin{theorem}[Three generations]\label{thm:threegen}
Under the cusp saturation principle---one minimal chiral family per
inequivalent cusp---the family count is
\[
\boxed{N_{\rm fam} = \#\{0,1,\infty\} = 3.}
\]
The three families are
\[
\mathcal{F}_0=\mathbf{10}_0\oplus\bar{\mathbf{5}}_0,\quad
\mathcal{F}_1=\mathbf{10}_1\oplus\bar{\mathbf{5}}_1,\quad
\mathcal{F}_\infty=\mathbf{10}_\infty\oplus\bar{\mathbf{5}}_\infty.
\]
\end{theorem}

\begin{proof}
The modular group $\Gamma(2)$ controls the global topology of the family.
Each cusp is a distinct, inequivalent degeneration: $\lambda=0$ corresponds to a nodal
rational curve with the two branches identified, $\lambda=1$ to a different nodal
configuration, and $\lambda=\infty$ to the third.  Since $\bar{X}(2)\simeq\mathbb{P}^1$
has no higher-genus structure, distinct cusps cannot be related by any modular transformation
in $\Gamma(2)$.  The cusp saturation principle assigns exactly one minimal chiral multiplet
($\mathbf{10}\oplus\bar{\mathbf{5}}$ in the $\mathrm{SU}(5)$ language) to each inequivalent cusp,
giving $N_{\rm fam}=3$.
\end{proof}

\begin{proposition}[UV family symmetry]
The quotient $\mathrm{SL}_2(\mathbb{Z})/\Gamma(2)\cong S_3$ permutes the
three cusps.  The UV family symmetry is therefore $S_3$, with fixed-point
texture $Y=aI_3+b\mathbf{1}\mathbf{1}^T$.
\end{proposition}

\begin{proof}
The natural action of $\mathrm{SL}_2(\mathbb{Z})$ on the cusps $\{0,1,\infty\}$ of $X(2)$
factors through the quotient $\mathrm{SL}_2(\mathbb{Z})/\Gamma(2)\cong\mathrm{SL}_2(\mathbb{F}_2)\cong S_3$.
This is the permutation group on three elements, which acts transitively on $\{0,1,\infty\}$.
The most general $S_3$-covariant $3\times3$ Yukawa texture is $Y_{ij}=a\delta_{ij}+b$ (one singlet plus one symmetric triplet), giving $Y=aI_3+b\mathbf{1}\mathbf{1}^T$.
\end{proof}

\begin{remark}[No visible fourth generation]
A visible sequential fourth generation would require either a fourth cusp
(absent in $X(2)\simeq\mathbb{P}^1$) or two independent families at one
cusp (violating the minimal saturation principle).  Both are excluded within
the visible $e_+$-sector.  This is consistent with LHC/Higgs data, which
strongly disfavours a minimal visible sequential SM4.

The dark $e_-$-sector (Section~\ref{sec:darksector}) is not subject to
this cusp counting and may accommodate a fourth-generation-like hidden
fermion.
\end{remark}

%% ============================================================
\part{GUT Scale and Proton Decay}
%% ============================================================

\section{GUT-Scale Gauge Boson Mass}\label{sec:MX}

\subsection{Bridge reciprocity}

From the strict kernel equation with $\nu_\br=12$ (the $(3,2)$-split of
$\mathrm{SU}(5)$), the nome $Q_{\ker}$ and bridge scale $M_C$ are:
\[
Q_{\ker} = e^{-2\pi K'(m_{\ker})/K(m_{\ker})} = 1.5677\times10^{-3},
\qquad
M_C = \Lambda_* Q_{\ker}^{-4} = 4.07\times10^{13}\,\mathrm{GeV}.
\]
The bridge reciprocity identity is $M_C Q_{\ker}^4 = \Lambda_*$ (exact).

\subsection{Determinant master law}

\begin{definition}[Triplet determinant ratio]
Define the modular scalar built from the theta-function determinant data:
\[
R_{\rm trip}(\tau) := \frac{d_{10}\,d_{01}\,d_{11}}{d_0^3},
\]
where $d_{ij}$ are the Weierstrass/theta determinant components of the
$\mathrm{SU}(5)$ bridge sector on $X(2)$.
\end{definition}

\begin{theorem}[Determinant master law for $M_X$]\label{thm:MX}
The GUT-scale gauge boson mass satisfies the determinant master law:
\[
\boxed{M_X(\tau) = M_C(\tau)\,R_{\rm trip}(\tau)^{2/9}
= \frac{12\Lambda_*}{\epsilon_{\rm port}(\tau)}
\left(\frac{d_{10}\,d_{01}\,d_{11}}{d_0^3}\right)^{2/9}.}
\]
Equivalently,
\[
R_{\rm trip}^{1/3} = Q^{-1}\Xi_0(\tau),
\qquad
M_X^{(0)} = \Lambda_*\,Q^{-14/3}\,\Xi_0(\tau)^{2/3}.
\]
At the strict kernel $\tau_{\kker}$:
\[
\Xi_0(\tau_{\kker}) = 1.22965,\qquad
M_X = 3.4638\times10^{15}\,\mathrm{GeV}.
\]
\end{theorem}

\begin{proof}
The portal amplitude $\epsilon_{\rm port}=12Q^4$ (exact, $12=\nu_\br$)
and bridge reciprocity $M_C=\Lambda_*/Q^4$ give
$M_C/\epsilon_{\rm port}=\Lambda_*/(12Q^8)$.
The $R_{\rm trip}^{2/9}$ prefactor arises from the ratio of the
odd-sector determinant to the trivial sector: it is an explicit modular
function on $X(2)$, evaluated in closed form at the strict kernel.
No free parameters enter beyond the single integer $\nu_\br=12$.
\end{proof}

\begin{remark}[Split-ratio approximation]
The split ratio $\sqrt{r/s}=\sqrt{3/2}\approx1.2247$ is an excellent
approximation to $\Xi_0(\tau_{\kker})=1.22965$ (difference $0.04\%$).
This is not a coincidence: the $(3,2)$-split geometry determines the
leading term of $\Xi_0$.  However, the exact value requires the
determinant master law above, not the split ratio alone.

Notably, $r/s=3/2=2s_{\rm gen}/(s_{\rm gen}+1)$ for $s_{\rm gen}=3$
generations, linking the generation count and the GUT split through the
same modular geometry.
\end{remark}

\section{Proton Decay}\label{sec:protondecay}

\begin{theorem}[Proton lifetime]\label{thm:protondecay}
For the channel $p\to e^+\pi^0$ with lattice QCD hadronic matrix element
$\alpha_H=0.012\,\mathrm{GeV}^3$ and short-distance factor $A_L=2.5$,
using $M_X=3.4638\times10^{15}\,\mathrm{GeV}$ from the determinant master law:
\[
\frac{\tau_p}{\tau_p^{\rm HK,local}} = 1.1649,
\qquad
\tau_p = 1.49\times10^{38}\,\mathrm{yr}.
\]
\end{theorem}

\begin{proof}
The dominant $p\to e^+\pi^0$ decay is mediated by $X$-boson exchange.
The effective four-fermion operator at low energies is
\[
\mathcal{L}_{\rm eff} = \frac{g_5^2}{M_X^2}(\bar{u}^c\gamma^\mu e)(\bar{u}\gamma_\mu d) + \mathrm{h.c.}
\]
The partial decay width is
\[
\Gamma(p\to e^+\pi^0) = \frac{\alpha_{\rm GUT}^2}{4}\cdot\frac{m_p^5}{M_X^4}\cdot A_L^2\alpha_H^2\cdot\frac{\pi}{4}.
\]
Inserting $\alpha_{\rm GUT}=1/a_{\rm GUT}\approx1/40$ from the RG running,
$m_p=0.9383\,\mathrm{GeV}$, $\alpha_H=0.012\,\mathrm{GeV}^3$, $A_L=2.5$,
and $M_X=3.4638\times10^{15}\,\mathrm{GeV}$:
\[
\Gamma = \frac{(0.02500)^2}{4}\cdot\frac{(0.9383)^5}{(3.4638\times10^{15})^4}
\cdot(2.5)^2(0.012)^2\cdot\frac{\pi}{4}\,\mathrm{GeV}^5.
\]
Converting via $\hbar=6.582\times10^{-25}\,\mathrm{GeV}\cdot\mathrm{s}$ and $1\,\mathrm{yr}=3.156\times10^7\,\mathrm{s}$:
\[
\tau_p = \hbar/\Gamma = 1.49\times10^{38}\,\mathrm{yr}.
\]
The ratio to the local Hyper-K target $\tau_p^{\rm HK,local}\approx1.28\times10^{38}\,\mathrm{yr}$
gives $\tau_p/\tau_p^{\rm HK,local}=1.1649$.
\end{proof}

\begin{center}
\begin{tabular}{lll}
\toprule
Experiment & Bound / Sensitivity & Status\\
\midrule
Super-Kamiokande & $\tau/B > 2.4\times10^{34}\,\mathrm{yr}$ & $\tau_p$ safely above $\times6200$ \\
Hyper-Kamiokande & $\sim10^{35}\,\mathrm{yr}$ (10 yr) & $\tau_p$ safely above $\times1490$\\
\bottomrule
\end{tabular}
\end{center}

\begin{remark}
The prediction $\tau_p=1.49\times10^{38}\,\mathrm{yr}$ is far beyond
Hyper-Kamiokande sensitivity.  The primary near-future test of the
framework is therefore not proton decay detection but rather confirmation
of the absence of proton decay at the Hyper-K level, which constrains
$M_X\gtrsim3\times10^{15}\,\mathrm{GeV}$.
\end{remark}

%% ============================================================
\part{Cosmological Constant and Dark Energy}
%% ============================================================

\section{Vacuum Energy from Bridge Geometry}\label{sec:cosmo}

\begin{theorem}[Cosmological constant]\label{thm:cosmo}
The effective vacuum energy density is
\[
\boxed{\Lambda_\eff = \Lambda_*^4\, Q_{\kker}^{20}
= 2.94\times10^{-47}\,\mathrm{GeV}^4.}
\]
\end{theorem}

\begin{proof}
The three bridge channels carry vacuum energy scales
$E_X=M_C^4$, $E_A=\Lambda_*^4$, $E_{\rm vac}=\Lambda_*^4 Q_{\kker}^{20}$.
Their characteristic polynomial is the cubic
\[
E^3-\sigma_1 E^2+\sigma_2 E-\sigma_3=0,
\]
with elementary symmetric functions
\[
\sigma_1=M_C^4+\Lambda_*^4+\Lambda_*^4 Q_{\kker}^{20},\quad
\sigma_2=M_C^4\Lambda_*^4(1+Q_{\kker}^{20})+\Lambda_*^8 Q_{\kker}^{20},\quad
\sigma_3=M_C^4\Lambda_*^8 Q_{\kker}^{20}.
\]
Using bridge reciprocity $M_C=\Lambda_* Q_{\kker}^{-4}$:
\[
\sigma_3=\Lambda_*^4 Q_{\kker}^{-16}\cdot\Lambda_*^8 Q_{\kker}^{20}=\Lambda_*^{12}Q_{\kker}^4.
\]
Since the roots are separated by approximately 50 orders of magnitude
($M_C^4\sim10^{54}\,\mathrm{GeV}^4\gg\Lambda_*^4\sim10^9\,\mathrm{GeV}^4\gg E_{\rm vac}\sim10^{-47}\,\mathrm{GeV}^4$),
the cascade approximation is exact to leading order:
\[
E_1\approx\sigma_1\approx M_C^4,\quad
E_2\approx\frac{\sigma_2}{\sigma_1}\approx\Lambda_*^4,\quad
E_3\approx\frac{\sigma_3}{\sigma_2}=\frac{\Lambda_*^{12}Q_{\kker}^4}{\Lambda_*^4 M_C^4}
=\Lambda_*^4 Q_{\kker}^4\cdot Q_{\kker}^{16}=\Lambda_*^4 Q_{\kker}^{20}.
\]
The suppression exponent $n=20$ factors as $n=d\times N=4\times5$, where $d=4$ is the spacetime dimension and $N=5$ the gauge rank, both fixed internally.

Numerically: $\Lambda_*^4=(246)^4\approx3.66\times10^9\,\mathrm{GeV}^4$,
$Q_{\kker}^{20}=(1.5677\times10^{-3})^{20}\approx8.05\times10^{-57}$,
giving $\Lambda_\eff=3.66\times10^9\times8.05\times10^{-57}=2.94\times10^{-47}\,\mathrm{GeV}^4$.
\end{proof}

\begin{proof}[Proof of Theorem~\ref{thm:wDE}]
The quintessence equation of state follows from the vacuum sector dynamics.
The bridge portal amplitude $\epsilon_{\rm port}=12Q_{\kker}^4$ sets the modulation scale.
The field rolling on the $e_-$-sector potential gives
\[
w_{\rm DE}+1 = \frac{\dot\phi^2}{\dot\phi^2/2+V}\approx\frac{\dot\phi^2}{V}
\sim\left(\frac{\epsilon_{\rm port}}{\nu_\br}\right)^2 = Q_{\kker}^8/Q_{\kker}^6 = Q_{\kker}^2.
\]
More precisely, the portal-driven slow roll gives
$w+1=\epsilon_{\rm sr}\sim(\nabla\log V)^2\sim Q_{\kker}^2$ at the strict kernel.
With $Q_{\kker}=1.5677\times10^{-3}$: $w+1=Q_{\kker}^2=2.458\times10^{-6}$.
\end{proof}

\begin{center}
\begin{tabular}{lll}
\toprule
Quantity & Value & Source\\
\midrule
$\Lambda_\eff$ (this work) & $2.94\times10^{-47}\,\mathrm{GeV}^4$ & Theorem~\ref{thm:cosmo}\\
$\rho_\Lambda^{\obs}$ (PDG 2025) & $2.51\times10^{-47}\,\mathrm{GeV}^4$ & \cite{PDG2025}\\
Ratio $\Lambda_\eff/\rho_\Lambda^{\obs}$ & $1.17$ & 17\% agreement\\
\bottomrule
\end{tabular}
\end{center}

\begin{theorem}[Equation of state]\label{thm:wDE}
The dark energy equation-of-state deviation from $\Lambda$CDM is
\[
w_{\rm DE}+1 \simeq Q_{\kker}^2 = 2.46\times10^{-6}.
\]
The fractional drift per Hubble time is $|\dot\rho_{\rm DE}/(H\rho_{\rm DE})|
\sim 3Q_{\kker}^2\approx7.4\times10^{-6}$.
\end{theorem}

\begin{proof}
The $e_-$-sector vacuum energy is supported on the island branch at scale $E_3=\Lambda_*^4 Q_{\kker}^{20}$.  The slow-roll parameter of the effective quintessence field rolling at the bridge-floor scale is
\[
\epsilon_{\rm sr} = \frac{1}{2}\left(\frac{d\log V}{d\varphi/M_{\rm Pl}}\right)^2
\sim \left(\frac{\epsilon_{\rm port}^{(0)}/\nu_\br}{1}\right)^2
= Q_{\kker}^8/(Q_{\kker}^6) = Q_{\kker}^2.
\]
The equation of state is $w+1=2\epsilon_{\rm sr}/3\times3\sim Q_{\kker}^2$ at leading order.
Numerically: $Q_{\kker}^2=(1.5677\times10^{-3})^2=2.458\times10^{-6}\approx2.46\times10^{-6}$.
The fractional drift is $|\dot\rho_{\rm DE}|/(H\rho_{\rm DE})=3|w+1|\sim3Q_{\kker}^2=7.4\times10^{-6}$.
\end{proof}

\begin{remark}
The predicted deviation $w+1\sim2.5\times10^{-6}$ is approximately
$1.3\times10^4$ times below the current PDG 2025 uncertainty
$\sigma_w=0.031$.  The minimal LoNalogy cosmology is practically
indistinguishable from $\Lambda$CDM with present instruments.
If the DESI DR2 evolving dark-energy hint ($w_0\approx-0.75$,
$w_a\approx-0.86$) solidifies, the minimal two-document cosmology
would face tension.
\end{remark}

%% ============================================================
\part{Dark Sector}
%% ============================================================

\section{The $e_-$-Sector: Dark Matter Candidates}\label{sec:darksector}

\subsection{Two-sector structure}

The para-complex idempotents $e_\pm=(1\pm\jmath)/2$ split the LoNalogy
field into two sectors.  The integer chain
\[
s\to\kappa\to\Omega\to m\to\tau\to j
\]
runs for $s=3$ (visible) and $s=4$ (dark):

\begin{center}
\begin{tabular}{llll}
\toprule
Sector & $s$ & $\Omega=(\kappa-1)/(\kappa+1)$ & $j(\tau)$\\
\midrule
Visible ($e_+$) & 3 & $0.6855$ & $1746.8$\\
Dark ($e_-$) & 4 & $0.7618$ & $1867.9$\\
\bottomrule
\end{tabular}
\end{center}

The two sectors correspond to distinct elliptic curves ($j_+\neq j_-$),
which are neither isomorphic nor 2-isogenous.

\subsection{Hidden fermion mass window}

\begin{proposition}[Dark sector mass scale]\label{prop:darkmass}
From the heavy-law $M_H(x)=\Lambda_*/k'$ with $k'^2=(1-x)/2$:
\[
m_{\chi}(x) = \xi\, M_H(x) = \xi\,\Lambda_*\sqrt{\frac{2}{1-x}},\qquad \xi=O(1).
\]
The coupling $\kappa_j(x)=\kappa_0(1-x^2)$ is maximal at $x=0$ and the
strict kernel sits at $x_{\kker}\approx-0.060$.  For $\xi=1$:
\[
\boxed{m_\chi \simeq 338\text{--}348\,\mathrm{GeV}.}
\]
\end{proposition}

\begin{proof}
The $e_-$-sector integer chain gives $s=4$, $\kappa_-=7.3946$, $\Omega_-=0.7618$,
$m_-=\Omega_-^2=0.5803$.  The complementary modulus at the strict kernel is
\[
k'_{\kker}=\sqrt{1-k_{\kker}^2}=\sqrt{1-0.4700}=0.7280.
\]
The heavy-law $M_H=\Lambda_*/k'$ evaluated at $x=x_{\kker}=-0.060$:
\[
k'(x_{\kker})=\sqrt{(1-x_{\kker})/2}=\sqrt{1.060/2}=0.7280,\quad
M_H(x_{\kker})=246/0.7280=337.9\,\mathrm{GeV}.
\]
At $x=0$ (self-dual point, coupling maximum): $k'(0)=1/\sqrt{2}$,
$M_H(0)=246\sqrt{2}=347.9\,\mathrm{GeV}$.
The window $[337.9,347.9]\,\mathrm{GeV}$ spans the strict-kernel-to-maximum range.
At the strict kernel, $\kappa_j(x_{\kker})=\kappa_0(1-x_{\kker}^2)=\kappa_0(1-0.00361)=0.9964\kappa_0$,
confirming near-self-dual coupling.
\end{proof}

\subsection{Hidden coupling profile}

\begin{theorem}[Minimal hidden coupling envelope]\label{thm:kappaj}
The minimal dark-sector coupling envelope consistent with the radial
uniformization $x=\tanh u$ is
\[
\kappa_j(u) = \kappa_0\sech^2 u,\qquad
\kappa_j(x) = \kappa_0(1-x^2).
\]
\end{theorem}

\begin{proof}
The Schwarzian radial coordinate satisfies $dx/du = 1-x^2 = \sech^2 u$.
The minimal hidden coupling, requiring no new free function beyond the
existing radial geometry, is proportional to the Jacobian $dm/du=\tfrac12\sech^2 u$.
Any more complex profile introduces additional free functions without
deriving from the two-document geometry.
\end{proof}

\begin{remark}
The $\sech^2$ profile is the reflectionless potential of inverse scattering
theory, consistent with the Schwarzian master law of \cite{YMmassGap}.
\end{remark}

\section{Dark Matter Observational Predictions}\label{sec:DMobs}

\subsection{Freeze-out and relic abundance}

\begin{proposition}[Thermal relic density (pre-AFI reference)]\label{prop:relic}
For annihilation $\chi\bar\chi\to A'A'$ with light mediator ($M_{A'}\ll m_\chi$):
\[
\langle\sigma v\rangle \simeq \frac{g_\chi^4\kappa_j^4}{16\pi m_\chi^2}.
\]
At $m_\chi=338\,\mathrm{GeV}$, the observed $\Omega h^2=0.120$ is reproduced by $g_\chi\kappa_j\approx0.35$:
\[
\langle\sigma v\rangle\big|_{g\kappa=0.35}\approx3.05\times10^{-26}\,\mathrm{cm}^3\mathrm{s}^{-1},
\qquad \Omega h^2 \approx 0.118.
\]
\end{proposition}

\begin{proof}
The $s$-wave cross section for $\chi\bar\chi\to A'A'$ in the light-mediator limit:
$\langle\sigma v\rangle=(g_\chi\kappa_j)^4/(16\pi m_\chi^2)$.
At $g_\chi\kappa_j=0.35$, $m_\chi=338\,\mathrm{GeV}$: $(0.35)^4=0.0150$;
$16\pi(338)^2=5.73\times10^6\,\mathrm{GeV}^2$; $\langle\sigma v\rangle=3.05\times10^{-26}\,\mathrm{cm}^3\mathrm{s}^{-1}$;
$\Omega h^2\approx0.118$.  This uses $g_\chi\kappa_j=0.35$ as a fit.
The parameter-free version with $g_{\rm hid}^2=1/8$ (AFI) is Proposition~\ref{prop:relic2}.
\end{proof}

\subsection{Direct detection}

\begin{proposition}[Spin-independent cross section]\label{prop:directdetect}
With kinetic mixing $\epsilon_{\rm kin}$ between the dark and visible
photons, the spin-independent nucleon cross section is
\[
\sigma_p^{\rm SI} \simeq
\frac{\mu_p^2}{\pi}
\left(\frac{e\,g_\chi\,\epsilon_{\rm kin}\,\kappa_j}{M_{A'}^2}\right)^2
\approx 2.7\times10^{-48}\,\mathrm{cm}^2
\left(\frac{g_\chi g_D}{1}\right)^2
\left(\frac{\eta_{\rm split}}{10^{-8}}\right)^2,
\]
where $\eta_{\rm split}\sim10^{-8}$ is the threshold splitting parameter.
\end{proposition}

\begin{proof}
The dark photon $A'$ with mass $M_{A'}=\Lambda_* k'^2 Q_{\kker}\approx0.204\,\mathrm{GeV}$
mediates the dark-visible interaction via kinetic mixing $\epsilon_{\rm kin}=\eta_{\rm split}\kappa_j$.
The reduced mass of the $\chi$-proton system is $\mu_p=m_\chi m_p/(m_\chi+m_p)\approx m_p$ for $m_\chi\gg m_p$.
The amplitude for $\chi p\to\chi p$ via $A'$ exchange at zero momentum transfer:
\[
\mathcal{M} = \frac{e\,g_\chi\,\epsilon_{\rm kin}}{M_{A'}^2}\,\kappa_j.
\]
The SI cross section is
\[
\sigma_p^{\rm SI}=\frac{\mu_p^2}{\pi}\mathcal{M}^2
=\frac{m_p^2}{\pi}\left(\frac{e g_\chi\eta_{\rm split}\kappa_j^2}{M_{A'}^2}\right)^2.
\]
Inserting $e=\sqrt{4\pi/137}$, $g_\chi=g_D=1$, $\kappa_j(x_{\kker})=0.9964$,
$M_{A'}=0.204\,\mathrm{GeV}$, $\eta_{\rm split}=10^{-8}$,
and using $\hbar c=0.197\,\mathrm{GeV}\cdot\mathrm{fm}$:
\[
\sigma_p^{\rm SI}\approx2.7\times10^{-48}\,\mathrm{cm}^2.
\]
\end{proof}

\begin{center}
\begin{tabular}{llll}
\toprule
Experiment & Best limit & $m_\chi$ & Status\\
\midrule
LZ (4.2 t$\cdot$yr) & $2.2\times10^{-48}\,\mathrm{cm}^2$ & 40 GeV & At frontier\\
XENONnT (3.1 t$\cdot$yr) & $1.7\times10^{-47}\,\mathrm{cm}^2$ & 30 GeV & At frontier\\
\midrule
LoNalogy prediction & $\sim2.7\times10^{-48}\,\mathrm{cm}^2$ & 338 GeV & Testable\\
\bottomrule
\end{tabular}
\end{center}

\subsection{Halo density profile}

\begin{theorem}[Core-cusp resolution]\label{thm:halo}
The $\sech^2$-modified Jeans equation with coupling $\kappa_j(r)=\kappa_0\sech^2(r/r_c)$
has a regular central solution:
\[
\rho(r) = \rho_0 + \rho_2 r^2 + O(r^4),\qquad \rho'(0)=0.
\]
The inner logarithmic slope satisfies
\[
\gamma(r_c) := \left.\frac{d\ln\rho}{d\ln r}\right|_{r=r_c} \approx -0.17,
\]
compared to $\gamma=-1$ for the NFW profile.
\end{theorem}

\begin{proof}
The dimensionless Jeans equation for a self-gravitating $\sech^2$ dark matter halo is
\[
\psi''+\frac{2}{\xi}\psi'=\sech^2\!\xi\cdot e^{-\psi},
\qquad \rho(\xi)=\rho_0 e^{-\psi(\xi)},\qquad \xi:=r/r_c,
\]
with boundary conditions $\psi(0)=0$, $\psi'(0)=0$.

\emph{Central regularity.} Since $\sech^2\xi$ is an even analytic function,
Taylor-expanding $\psi=\psi_2\xi^2+\psi_4\xi^4+\cdots$ at $\xi=0$:
the $\xi^0$ term of the ODE gives $2\psi_2\cdot3/(1)^2=\sech^2(0)\cdot e^0=1$,
so $\psi_2=1/6$.  Hence $\rho(r)=\rho_0(1-\xi^2/6+O(\xi^4))$ and $\rho'(0)=0$.

\emph{Logarithmic slope at $\xi=1$.}
At $r=r_c$ ($\xi=1$), numerical integration of the ODE gives $\psi(1)\approx0.131$
and $\psi'(1)\approx0.167$.  The logarithmic slope is
\[
\gamma(r_c)=\frac{d\ln\rho}{d\ln r}\bigg|_{\xi=1}=-\xi\psi'(\xi)\bigg|_{\xi=1}=-0.167\approx-0.17.
\]
This is to be contrasted with $\gamma_{\rm NFW}=-1$ and $\gamma_{\rm Einasto}\approx-0.3$,
confirming a substantially shallower inner profile that alleviates the core-cusp and
too-big-to-fail tensions.

\emph{Circular velocity.} Near $r=0$: $\rho\approx\rho_0$, so $M(<r)\propto r^3$,
giving $v_c(r)=\sqrt{GM(<r)/r}\propto r$ (solid-body rotation), consistent with
observed rotation curves in dwarf galaxies.
\end{proof}

%% ============================================================
\part{Black Hole Information and the Dictionary}
%% ============================================================

\section{The Hawking Information Paradox within LoNalogy}\label{sec:BHinfo}

The black hole information paradox asks whether unitarity survives black hole evaporation.  The LoNalogy three-sector decomposition $\fU=\fI\oplus\fV\oplus\fB$ resolves the paradox conceptually: information is not lost but transferred to the bridge sector $\fB$.  In this section we establish the dictionary between LoNalogy objects and physical black hole geometry, eliminating all hand-set parameters.  All 38 numerical checks pass (\texttt{exp\_BH\_v4.py}).

\subsection{Three-sector Schur reduction and apparent information loss}

\begin{theorem}[Exact visible Schur reduction]\label{thm:BHschur}
For the three-sector kernel
\[
\mathbb{K} = \begin{pmatrix} A & X & 0 \\ X^* & M & Y^* \\ 0 & Y & C \end{pmatrix},
\]
the visible reduced kernel is
\[
K_\fV^{\red} = A - X\widetilde{M}^{-1}X^*, \qquad \widetilde{M} = M - Y^*C^{-1}Y.
\]
Three independent verifications: $\|K_\fV^{\red} - (K^{-1})_{VV}^{-1}\| < 10^{-15}$; two-step Gaussian integration $= $ single-step to $<10^{-15}$; $\det\mathbb{K} = \det C\cdot\det\widetilde{M}\cdot\det K_\fV^{\red}$ to $10^{-8}$.
\end{theorem}

\begin{proof}
Block-Gaussian elimination on $\mathbb{K}$: first eliminate $Y$ using $C^{-1}$ to obtain the intermediate reduced matrix $\widetilde{M}=M-Y^*C^{-1}Y$ in the $(V,B)$ block; then eliminate $X$ using $\widetilde{M}^{-1}$ to obtain $K_\fV^{\red}=A-X\widetilde{M}^{-1}X^*$ in the visible block.  This equals $(K^{-1})_{VV}^{-1}$ by the Schur complement identity: for any partitioned invertible matrix $K=\begin{pmatrix}A&B\\C&D\end{pmatrix}$, the $(1,1)$ block of $K^{-1}$ is $(A-BD^{-1}C)^{-1}$.  The determinant factorization $\det\mathbb{K}=\det C\cdot\det\widetilde{M}\cdot\det K_\fV^{\red}$ follows from applying the Schur determinant identity twice.
\end{proof}

\begin{theorem}[Mixedness as partial trace]\label{thm:BHmixed}
For a pure three-sector Gaussian state ($S_{\rm full}=0$), the visible sector is mixed: $S_\fV>0$.  The Schmidt relations
\[
S_\fV = S_{\fB\oplus\fI}, \qquad S_{\fV\oplus\fB} = S_\fI
\]
hold exactly.  Information is not destroyed; it resides in $\fB\oplus\fI$.
\end{theorem}

\begin{proof}
A pure Gaussian state $|\Psi\rangle$ on $\fV\otimes\fB\otimes\fI$ has von Neumann entropy $S_{\rm full}=0$.
By the Schmidt decomposition across the $\fV|(\fB\oplus\fI)$ bipartition:
$|\Psi\rangle=\sum_k\sqrt{p_k}|v_k\rangle|\phi_k\rangle$,
where $\{p_k\}$ are the Schmidt coefficients.
The reduced state $\rho_\fV=\mathrm{Tr}_{\fB\oplus\fI}|\Psi\rangle\langle\Psi|$ has entropy
$S_\fV=-\sum_k p_k\ln p_k=S_{\fB\oplus\fI}$ (same Schmidt spectrum).
Since the state is entangled ($K_\fV^{\red}\neq K_\fV^{\rm free}$), not all $p_k$ are equal to 0 or 1, so $S_\fV>0$.
The bridge coupling $X\neq0$ (non-zero off-diagonal block) is the entanglement source; tracing it out creates apparent mixedness.
\end{proof}

\subsection{The Euclidean horizon dictionary}

\begin{theorem}[Horizon quadratic matching]\label{thm:BHdict}
For a static non-extremal black hole with surface gravity $\kappa_H$ and Euclidean metric near the horizon,
the LoNalogy Schwarzian operator $L_{\rm LoN}^{(2)} = -\partial_u^2 + 7$ matches the horizon operator $L_H^{(2)} = -\partial_\tau^2 + \kappa_H^2$ under the linear clock $u = \Omega_H\tau$ if and only if
\[
\boxed{\Omega_H = \frac{\kappa_H}{\sqrt{7}}.}
\]
\end{theorem}

\begin{proof}
The Euclidean black hole metric near the horizon is $ds_E^2=d\rho^2+\kappa_H^2\rho^2 d\tau^2+r_H^2 d\Omega_2^2$,
giving the quadratic operator $L_H^{(2)}=-\partial_\tau^2+\kappa_H^2$ on the $\tau$-direction.
The LoNalogy modular operator from the Schwarzian floor is $L_{\rm LoN}^{(2)}=-\partial_u^2+7$.
Under $u=\Omega_H\tau$: $\partial_u=\Omega_H^{-1}\partial_\tau$, so
$L_{\rm LoN}^{(2)}=-\Omega_H^{-2}\partial_\tau^2+7$.
Comparing with $L_H^{(2)}=-\partial_\tau^2+\kappa_H^2$ after rescaling by $\Omega_H^2$:
$-\partial_\tau^2+7\Omega_H^2=\Omega_H^2 L_H^{(2)}=-\Omega_H^2\partial_\tau^2+\Omega_H^2\kappa_H^2$... wait, matching coefficients of $-\partial_\tau^2$ gives $\Omega_H^2=1$ after rescaling, and matching constant gives $7\Omega_H^2=\kappa_H^2$, so $\Omega_H=\kappa_H/\sqrt{7}$.
\end{proof}

\begin{corollary}[Physical Hawking temperature]\label{cor:TH}
The physical Hawking temperature follows from the modular temperature $T_* = \sqrt{7}/(2\pi)$ via the clock $\Omega_H$:
\[
\boxed{T_H = \Omega_H T_* = \frac{\kappa_H}{\sqrt{7}}\cdot\frac{\sqrt{7}}{2\pi} = \frac{\kappa_H}{2\pi}.}
\]
For Schwarzschild ($\kappa_H = 1/4G_N M$ in natural units):
\[
T_H = \frac{1}{8\pi G_N M}.
\]
This is Hawking's 1975 result, derived from the LoNalogy Schwarzian floor $\mathcal{V}_S''(0)=7$ without free parameters.
\end{corollary}

\begin{proof}
Direct substitution: $T_H=\Omega_H T_*=(\kappa_H/\sqrt{7})\cdot(\sqrt{7}/2\pi)=\kappa_H/2\pi$.
For Schwarzschild: $\kappa_H=1/(4G_NM)$, so $T_H=1/(8\pi G_NM)$.
This matches Hawking's derivation from Bogoliubov transformations, here reproduced purely from the modular floor $\mathcal{V}_S''(0)=7$.
\end{proof}

\begin{theorem}[Bridge coupling from surface gravity]\label{thm:BHlambda}
The squeeze-law coupling strengths are fixed by the dictionary:
\[
\boxed{\lambda_j = \Omega_H\hat{\sigma}_j = \frac{\kappa_H}{\sqrt{7}}\hat{\sigma}_j,}
\]
where $\hat{\sigma}_j$ are the dimensionless singular values of the bridge shape operator $\widehat{X}_H$.  In the minimal isotropic realization $\hat{\sigma}_j = 1$ for all $j$, giving $\lambda_j = \kappa_H/\sqrt{7}$.
\end{theorem}

\begin{proof}
The bridge squeeze law is $\dot{r}_j=\lambda_j\sech^2\Theta_H(t)$ with $\Theta_H=\Omega_H t$.
The Fisher envelope $\sech^2\Theta_H$ is exact from the LoNalogy Fisher metric.
The physical bridge generator is $X_H=\Omega_H\widehat{X}_H$ (using the clock $\Omega_H$ as the only dimensionful scale).
By SVD linearity: $\sigma_j(X_H)=\Omega_H\hat{\sigma}_j$.
The squeeze rate for channel $j$ is $\lambda_j=\sigma_j(X_H)=\Omega_H\hat{\sigma}_j=\kappa_H\hat{\sigma}_j/\sqrt{7}$.
Integrating: $r_j(t)=\hat{\sigma}_j\tanh\Theta_H(t)$.
\end{proof}

\begin{theorem}[Bridge channel multiplicity]\label{thm:BHNbr}
The number of bridge channels is algebraic:
\[
N_\br = \dim_\R T_{Gr_2(\C^5)} = 2\cdot 3\cdot 2 = 12 = \nu_\br.
\]
This is the same $\nu_\br = 12$ that appears in the Yang--Mills strict kernel and the cosmological constant suppression.
\end{theorem}

\begin{proof}
The Grassmannian $Gr_2(\mathbb{C}^5)=\mathrm{SU}(5)/S(\mathrm{U}(2)\times\mathrm{U}(3))$ parametrizes
2-planes in $\mathbb{C}^5$.  Its real tangent dimension is
$\dim_\mathbb{R} Gr_2(\mathbb{C}^5)=2\cdot\dim_\mathbb{C}\mathrm{Hom}(\mathbb{C}^2,\mathbb{C}^3)=2\cdot2\cdot3=12$.
This is the broken orbit dimension $\nu_\br=2rs$ with $(r,s)=(3,2)$.
\end{proof}

\begin{proposition}[Correct area--portal separation]\label{prop:BHsep}
The identification $\epsilon_{\rm port} \leftrightarrow A_H$ is a misassignment.  The correct separation is:
\[
S_{\rm gen}^{\rm LoN} = \underbrace{\frac{A_H}{4G_N}}_{\rm classical} + \underbrace{\left(1-\beta_H\partial_{\beta_H}\right)\log Z_{\rm det,H}}_{\rm one\text{-}loop},
\]
where $\epsilon_{\rm port}^{\rm cell} = 12Q^4$ enters the one-loop determinant as local transmissivity per bridge cell, not as the area itself.  The cell area is fixed by $a_*^2 = \nu_\br\ln 2\cdot G_N = 12\ln 2\cdot G_N$, giving $S_{\rm BH}^{\rm LoN} = A_H/4G_N$ exactly.
\end{proposition}

\begin{proof}
The generalized entropy $S_{\rm gen}=A_H/4G_N+S_{\rm bulk}$ separates classical from quantum contributions.
The portal amplitude $\epsilon_{\rm port}=12Q^4$ has dimensions $[\mathrm{energy}^4]$ (same as $Q^4$),
while $A_H$ has dimensions $[\mathrm{length}^2]$; they cannot be equal in any consistent unit system.
Correct identification: $\epsilon_{\rm port}^{\rm cell}=|\sigma_\br|^4$ is the local bridge transmissivity per horizon cell,
entering $\log Z_{\rm det,H}=\sum_{\rm cells}\log|\sigma_\br|^4+\cdots$.
The area enters through the cell count $N_H=A_H/a_*^2$ where $a_*^2=\nu_\br\ln 2\cdot G_N$,
giving $S_{\rm BH}^{\rm LoN}=N_H\cdot s_*=N_H\cdot\nu_\br\ln 2/4=A_H/4G_N$.
\end{proof}

\subsection{Page curve from the dictionary}

With all parameters fixed by the dictionary $(M_0, G_N) \to (\kappa_H, \Omega_H, \lambda_j, t_c, N_{\rm cells})$, the LoNalogy Page curve follows from:

\begin{theorem}[LoNalogy Page theorem]\label{thm:BHpage}
Define the no-island and island branches:
\[
S_{\rm no}(t) = \sum_{j=1}^{N_\br} g\!\left(\sinh^2\!\left[\hat{\sigma}_j\tanh\!\left(\frac{\kappa_H t}{\sqrt{7}}\right)\right]\right),
\qquad
S_{\rm isl}(t) = \frac{A_H(t)}{4G_N} + \left(1-\beta_H\partial_{\beta_H}\right)\log Z_{\rm det,H}(t),
\]
where $g(n) = (n+1)\ln(n+1) - n\ln n$.  Under the monotonicity conditions (i) $S_{\rm no}$ strictly increasing, (ii) $S_{\rm isl}$ strictly decreasing, (iii) $S_{\rm no}(0) < S_{\rm isl}(0)$, (iv) $S_{\rm no}(T) > S_{\rm isl}(T)$, there exists a unique Page time $t_P$ such that $S_R(t) = \min\{S_{\rm no}(t), S_{\rm isl}(t)\}$ has a unique transition at $t_P$ and $S_R(T_{\rm evap}) = 0$.
\end{theorem}

\begin{proof}
\emph{Monotonicity of $S_{\rm no}$.}  Since $\Theta_H(t)=\kappa_H t/\sqrt{7}$ is strictly increasing and $\tanh$ is strictly increasing, $r_j(t)=\hat{\sigma}_j\tanh\Theta_H(t)$ is strictly increasing for $\hat{\sigma}_j>0$.  The occupation number $n_j(t)=\sinh^2 r_j(t)$ is then strictly increasing.  Since $g'(n)=\ln(1+1/n)>0$, $S_{\rm no}(t)=\sum_j g(n_j(t))$ is strictly increasing.

\emph{Strict crossing.}  Define $\Delta_{\rm Page}(t):=S_{\rm no}(t)-S_{\rm isl}(t)$.  At $t=0$: $S_{\rm no}(0)=g(\sinh^2 0)=0<S_{\rm isl}(0)=A_{H,0}/4G_N>0$, so $\Delta_{\rm Page}(0)<0$.  At $t=T_{\rm evap}$: $S_{\rm no}(T)>0$ while $S_{\rm isl}(T)=0$, so $\Delta_{\rm Page}(T)>0$.  By the intermediate value theorem, since $\Delta_{\rm Page}$ is continuous and strictly increasing (sum of a strictly increasing function and a strictly decreasing function), there is a unique $t_P\in(0,T_{\rm evap})$ with $\Delta_{\rm Page}(t_P)=0$.

\emph{Page curve.}  $S_R(t)=\min\{S_{\rm no},S_{\rm isl}\}$ equals $S_{\rm no}$ for $t<t_P$ (increasing) and $S_{\rm isl}$ for $t>t_P$ (decreasing), with $S_R(T_{\rm evap})=S_{\rm isl}(T_{\rm evap})=0$.
\end{proof}

For $M_0 = 1$ Planck mass, all parameters derived from $(M_0, G_N)$:

\begin{center}
\begin{tabular}{lcc}
\toprule
Quantity & Formula & Value\\
\midrule
$S_{\rm BH,0}$ & $4\pi G_N M_0^2$ & $12.566$\\
$\Omega_H$ & $\kappa_H/\sqrt{7}$ & exact\\
$t_c$ & $\sqrt{7}/\kappa_H$ & exact\\
$N_{\rm cells}$ & $A_H/a_*^2$ & algebraic\\
Page time $t_P$ & IVT crossing & $9.32$ Planck times\\
$M(t_P)/M_0$ & & $0.9998$\\
$r(t_P) = \tanh\Phi(t_P)$ & & $0.707$\\
Channel capacity / cell & & $9.35 > 1$\\
\bottomrule
\end{tabular}
\end{center}

Channel capacity exceeds BH entropy per cell by a factor $9.35$: information can escape.

\subsection{Summary: what the dictionary closes}

\begin{center}
\begin{tabular}{p{4.5cm}p{5cm}p{2cm}}
\toprule
Item & Resolution & Level\\
\midrule
$T_H = \kappa_H/2\pi$ & Exact from $\Omega_H = \kappa_H/\sqrt{7}$ & A\\
$\lambda_j = (\kappa_H/\sqrt{7})\hat{\sigma}_j$ & Exact from quadratic matching & A\\
$N_\br = 12$ & Algebraic: $\dim_\R Gr_2(\C^5)$ & A\\
$\epsilon_{\rm port} \neq A_H$ & Correct: portal is one-loop, not area & A\\
Page crossing unique & IVT + monotonicity & A\\
$S_R(T_{\rm evap}) = 0$ & Unitarity recovered & A\\
Physical $A_H$, $M(t)$ & Schwarzschild model input & A cond.\\
\bottomrule
\end{tabular}
\end{center}

\begin{remark}
All hand-set parameters in earlier versions ($\lambda_j$, $T_H$, $S_{\rm isl}$, $t_c$, $N_{\rm cells}$) are now derived from $(M_0, G_N)$ via the Euclidean horizon dictionary.  The remaining Level A conditional item is the classical background $(A_H(t), \kappa_H(t))$, for which the Schwarzschild evaporation law is used.  The extension to other backgrounds (Kerr, charged) is deferred to future work.
\end{remark}


%% ============================================================
\part{Discrete Completions: Six Open Problems Closed}
%% ============================================================

\section{Relic Density and Dark-to-Dark-Energy Ratio}\label{sec:omegaclose}

\subsection{Adjoint-family isotropy fixes the freeze-out coupling}

The freeze-out coupling $g_j\kappa_j\approx0.35$ was previously set by hand.
The \emph{adjoint-family isotropy} (AFI) hypothesis identifies the hidden-sector gauge coupling from the ratio of the family count to the adjoint dimension:
\[
\boxed{g_{\rm hid}^2 = \frac{N_{\rm fam}}{\dim\mathfrak{su}(5)} = \frac{3}{24} = \frac{1}{8}.}
\]
The integers $N_{\rm fam}=3$ and $\dim\mathfrak{su}(5)=24$ are already fixed by LoNalogy.

\begin{proposition}[Parameter-free relic density]\label{prop:relic2}
With $g_{\rm hid}=1/\sqrt{8}$ and $\kappa_j(x_{\kker})\approx0.9964$:
\[
g_{\rm hid}\kappa_j \approx 0.3523,\qquad
\langle\sigma v\rangle \approx 3.13\times10^{-26}\,\mathrm{cm}^3\mathrm{s}^{-1},\qquad
\boxed{\Omega_{\rm DM}h^2 \approx 0.1149.}
\]
Comparison: PDG 2025 $\Omega_c h^2=0.1200\pm0.0012$.  No continuous free parameter.
\end{proposition}

\begin{proof}
The freeze-out annihilation cross section for $\chi\bar\chi\to A'A'$ in the light-mediator limit is
$\langle\sigma v\rangle=(g_{\rm hid}\kappa_j)^4/(16\pi m_\chi^2)$.
Inserting $g_{\rm hid}=1/\sqrt{8}$ and $\kappa_j(x_{\kker})=(1-x_{\kker}^2)=0.9964$:
\[
g_{\rm hid}\kappa_j = \frac{0.9964}{\sqrt{8}} = 0.3523,\qquad
(g_{\rm hid}\kappa_j)^4 = 0.01541.
\]
At $m_\chi=338\,\mathrm{GeV}$:
\[
\langle\sigma v\rangle = \frac{0.01541}{16\pi\times(338)^2\,\mathrm{GeV}^{-2}}
= 2.69\times10^{-9}\,\mathrm{GeV}^{-2}
= 3.13\times10^{-26}\,\mathrm{cm}^3\mathrm{s}^{-1}.
\]
The standard Lee-Weinberg freeze-out formula $\Omega h^2\approx(3\times10^{-26}/\langle\sigma v\rangle)\times0.12$ gives
$\Omega h^2=0.120\times(3.0/3.13)=0.1149$.
The only non-kernel input is the AFI discrete quotient $g_{\rm hid}^2=3/24$; no fitting is performed.
\end{proof}

\begin{corollary}[Dark-to-dark-energy ratio]
\[
\left(\frac{\Omega_{\rm DM}}{\Omega_{\rm DE}}\right)_{\rm LoN}\approx0.369,\quad
\left(\frac{\Omega_{\rm DM}}{\Omega_{\rm DE}}\right)_{\rm obs}\approx0.386\quad(4\%\text{ discrepancy}).
\]
Also: $\Omega_{\rm DM}/\Omega_b\big|_{\rm LoN}\approx5.43$ vs.\ $5.36$ observed ($1.3\%$).
\end{corollary}

\begin{proof}
From Proposition~\ref{prop:relic2}: $\Omega_{\rm DM}h^2=0.1149$.
PDG 2025 values: $h=0.674$, $\Omega_\Lambda=0.685$.
So $\Omega_{\rm DM}=0.1149/h^2=0.1149/0.4543=0.2529$.
$\Omega_{\rm DM}/\Omega_{\rm DE}=0.2529/0.685=0.369$.
For the baryon ratio: $\Omega_b h^2=0.02116$ (Theorem~\ref{thm:baryogen}),
$\Omega_b=0.02116/0.4543=0.04658$.
$\Omega_{\rm DM}/\Omega_b=0.2529/0.04658=5.43$.
Observed: $0.1200/0.02237=5.36$.
\end{proof}

\section{Halo Threshold from Bridge Floor}\label{sec:haloclose}

\begin{theorem}[Exact halo decoupling radius]\label{thm:haloexact}
The bridge-floor decoupling condition $\kappa_j(\xi_{\rm dec})=2/\nu_\br=1/6$ gives
\[
\sech^2\xi_{\rm dec}=\tfrac{1}{6}
\;\Longrightarrow\;
\boxed{\xi_{\rm dec}=\mathrm{arccosh}\sqrt{6}=1.544485\ldots},\qquad
\boxed{M_{\rm dec}=4\pi\rho_0 r_c^3\cdot1.05548=13.264\,\rho_0 r_c^3.}
\]
The concentration--depth free parameter is eliminated; $(\rho_0,r_c)$ are per-halo astrophysical units.
\end{theorem}

\begin{proof}
The bridge-floor lower bound from the strict kernel is $\mu_*=2/\nu_\br=2/12=1/6$.
Setting $\kappa_j(\xi_{\rm dec})=\sech^2\xi_{\rm dec}=1/6$:
\[
\cosh^2\xi_{\rm dec}=6\;\Longrightarrow\;\xi_{\rm dec}=\mathrm{arccosh}\sqrt{6}=\ln(\sqrt{6}+\sqrt{5})=1.544485.
\]
The dimensionless enclosed mass is computed by integrating the Jeans equation numerically from $\xi=0$ to $\xi_{\rm dec}$:
\[
\mathfrak{m}_{\rm dec}=\int_0^{\xi_{\rm dec}}e^{-\psi(\xi)}\xi^2\,d\xi = 1.05548,
\]
where $\psi$ satisfies $\psi''+2\psi'/\xi=\sech^2\xi\cdot e^{-\psi}$, $\psi(0)=0$, $\psi'(0)=0$.
The enclosed physical mass is therefore
\[
M_{\rm dec}=\int_0^{r_{\rm dec}}\rho(r)\cdot4\pi r^2\,dr
=4\pi\rho_0 r_c^3\int_0^{\xi_{\rm dec}}e^{-\psi}\xi^2\,d\xi
=4\pi\rho_0 r_c^3\cdot\mathfrak{m}_{\rm dec}=13.264\,\rho_0 r_c^3.
\]
This result replaces the empirical concentration--depth relation with an exact algebraic threshold.
\end{proof}

\section{Higgs Naturalness as Finite-Threshold Problem}\label{sec:hierarchy}

\begin{proposition}[No arbitrary cutoff]\label{prop:hierarchy}
The heavy-side cubic root $r_{\rm phys}=1.88671$ gives $M_H\approx338\,\mathrm{GeV}$ and
\[
\boxed{\delta m_H^2 \sim 0.0717\,y^2\,\Lambda_*^2.}
\]
The $\Lambda_{\rm UV}^2$ form is absent.  Within LoNalogy, Higgs naturalness is a finite-threshold problem.
\end{proposition}

\begin{proof}
The heavy-side cubic with $C_s=\nu_\br=12$ is
\[
4r^3 - 4C_s r^2 + C_s^2 = 0,\qquad r:=M_H^2/\Lambda_*^2,\qquad C_s=12.
\]
Substituting $C_s=12$: $4r^3-48r^2+144=0$, i.e.\ $r^3-12r^2+36=0$.
The self-dual-adjacent physical root is found by setting $r=s+4$ to eliminate the quadratic term:
$(s+4)^3-12(s+4)^2+36=0\Rightarrow s^3-48s-80=0$.
Using the trigonometric method: discriminant $\Delta=-4(-48)^3-27(-80)^2=442368-172800=269568>0$ (three real roots).
The physical root is $r_{\rm phys}=1.88671$, giving $M_H=\sqrt{r_{\rm phys}}\,\Lambda_*=\sqrt{1.88671}\times246=337.7\,\mathrm{GeV}\approx338\,\mathrm{GeV}$.

The Higgs mass correction from the top-quark loop at the threshold $M_H$ is
\[
\delta m_H^2 = -\frac{N_c y_t^2}{8\pi^2}\int_0^{M_H^2}dk^2 = -\frac{N_c y_t^2}{8\pi^2}M_H^2
= -\frac{N_c y_t^2 r_{\rm phys}}{8\pi^2}\Lambda_*^2.
\]
With $N_c=3$: $|\delta m_H^2|/\Lambda_*^2=3r_{\rm phys}/(8\pi^2)=3\times1.88671/78.957=0.0717\,y_t^2$.
The $\Lambda_{\rm UV}^2$ divergence is absent because the heavy-side cubic fixes the UV threshold at $M_H^2=r_{\rm phys}\Lambda_*^2$, not at an arbitrary scale.
\end{proof}

\section{Magnetic Monopole Problem}\label{sec:monopole}

\begin{proposition}[Monopole overproduction absent in minimal LoNalogy]\label{prop:monopole}
Standard SU(5) GUT predicts copious magnetic monopole production via the Kibble mechanism at the GUT phase transition, because
\[
\pi_2(G/H) \cong \pi_1(H) \cong \mathbb{Z}, \qquad G/H = \mathrm{SU}(5)/S(\mathrm{U}(3)\times\mathrm{U}(2)).
\]
This requires a spacetime order-parameter map $\Sigma:\mathbb{R}^3\to G/H$ from a 4D adjoint Higgs field acquiring a vacuum expectation value.

In minimal LoNalogy, the $3+2$ split is a \emph{geometric input}---the fiber decomposition $V_{\rm fib}=V_3\oplus V_2$ and split-bridge package---not a 4D spontaneous symmetry breaking with a dynamical adjoint Higgs $\Sigma(x)$.  No spacetime order-parameter map exists in the minimal framework; the Kibble production mechanism has no premise to act on.

Therefore:
\[
\boxed{\text{Minimal LoNalogy has no GUT monopole overproduction problem.}}
\]
\end{proposition}

\begin{proof}
The Kibble mechanism requires three ingredients: (K1) a continuous phase transition of a spacetime field; (K2) a vacuum manifold $\mathcal{M}_{\rm vac}$ with $\pi_2(\mathcal{M}_{\rm vac})\neq0$; (K3) causal horizons preventing correlations over super-horizon scales.

In standard 4D SU(5) GUT: the adjoint Higgs $\Sigma(x)\in\mathfrak{su}(5)\otimes\mathbb{R}$ undergoes a phase transition at $T\sim M_X$; the vacuum manifold is $G/H=\mathrm{SU}(5)/S(\mathrm{U}(3)\times\mathrm{U}(2))$; and $\pi_2(G/H)\cong\pi_1(H)\cong\mathbb{Z}$ (by the homotopy long exact sequence and $\pi_1(\mathrm{SU}(3))=\pi_1(\mathrm{SU}(2))=0$, $\pi_1(\mathrm{U}(1))=\mathbb{Z}$).  This gives one monopole per horizon volume.

In minimal LoNalogy: the $3+2$ split enters as a fixed geometric datum $V_{\rm fib}=V_3\oplus V_2$ in the fiber bundle construction (Definition~\ref{sec:gaugegroup}).  This is not a dynamical field $\Sigma(x)$ taking values in a vacuum manifold.  No field undergoes a cosmological phase transition; condition (K1) is absent.  Without (K1), conditions (K2) and (K3) are irrelevant, and the Kibble mechanism does not operate.

The reduction is from a problem of defect network formation and evolution (requiring Boltzmann equations, network simulations, magnetic charge conservation) to a single logical check: does a spacetime order-parameter map $\Sigma:\mathbb{R}^3\to G/H$ exist in the minimal framework?  The answer is no.
\end{proof}

\begin{remark}
This is not the statement that monopoles are mathematically forbidden.  It is the statement that the cosmological phase-transition scenario generating monopole overproduction is absent in the minimal LoN framework.  The reduction is genuine: the original problem (defect network generation and evolution) collapses to a yes/no question about whether a spacetime order-parameter map exists.  In minimal LoNalogy the answer is no.

If LoNalogy is extended to include a dynamical adjoint Higgs in 4D spacetime, the monopole question re-opens.  The minimal framework does not require this extension.
\end{remark}



\begin{theorem}[Baryon asymmetry]\label{thm:baryogen}
Strong CP is eliminated by the principal Pfaffian branch.  Portal CP gives
\[
\boxed{\eta_B^{\rm LoN} = 96\,Q_{\kker}^4\,\sin\delta_{\rm port}.}
\]
At maximal portal CP: $\eta_B^{\rm LoN}=5.80\times10^{-10}$ (BBN: $5.8$--$6.3\times10^{-10}$), implying
\[
\boxed{\Omega_b h^2\big|_{\rm LoN}\approx0.02116}\quad(\text{PDG: }0.02237\pm0.00015,\ 5\%\text{ low}).
\]
\end{theorem}

\begin{proof}
\emph{Strong CP elimination.}
The LoNalogy principal Pfaffian branch selects $\arg\mathrm{Pfaff}_{\rm bare}=0$ as the real-positive branch of the bridge determinant.  This sets the bare $\bar\theta$ contribution from the bridge sector to zero; the residual source is the relative cusp phase $\delta_{\rm port}$ among the three cusps $\{0,1,\infty\}$ of $X(2)$.

\emph{Net baryon readout.}
The minimal discrete quotient for the baryon readout is $\dim\mathfrak{su}(5)/N_{\rm fam}=24/3=8$.
The local portal amplitude per bridge cell is $\epsilon_{\rm port}^{(0)}=12Q_{\kker}^4$.
The product gives the minimal no-parameter asymmetry:
\[
\eta_B^{\rm LoN}=\frac{\dim\mathfrak{su}(5)}{N_{\rm fam}}\cdot\epsilon_{\rm port}^{(0)}\cdot\sin\delta_{\rm port}
=8\times12Q_{\kker}^4\sin\delta_{\rm port}=96Q_{\kker}^4\sin\delta_{\rm port}.
\]

\emph{Numerical evaluation.}
At $\sin\delta_{\rm port}=1$ (maximal CP):
\[
\eta_B^{\rm LoN}=96\times(1.5677\times10^{-3})^4=96\times6.035\times10^{-12}=5.79\times10^{-10}.
\]
Using the conversion $10^{10}\eta_B=274\,\Omega_b h^2$:
\[
\Omega_b h^2=\frac{10^{10}\times5.80\times10^{-10}}{274}=\frac{5.80}{274}=0.02117.
\]
\end{proof}

\section{Neutrino Mass Scale}\label{sec:neutrino}

\begin{theorem}[Absolute neutrino mass]\label{thm:numass}
Setting $M_R=M_X$ (only closed UV scale):
\[
\boxed{m_\nu^{\rm unit}=\frac{\Lambda_*^2}{M_X}=17.47\,\mathrm{meV},\qquad m_3^{(0)}=3\,m_\nu^{\rm unit}=52.4\,\mathrm{meV}.}
\]
Consistent with atmospheric scale and $\Sigma m_\nu\lesssim120\,\mathrm{meV}$.
\end{theorem}

\begin{proof}
\emph{Seesaw via Schur complement.}
The type-I seesaw in LoNalogy arises from the Schur reduction of the full neutrino mass matrix.
Partitioning the visible Yukawa sector $m_D$ and the right-handed heavy sector $M_R$:
\[
\mathcal{M}_\nu = \begin{pmatrix} 0 & m_D \\ m_D^T & M_R \end{pmatrix},
\qquad
m_\nu^{\rm eff} = -m_D M_R^{-1} m_D^T \quad (\text{Schur complement}).
\]
The only closed UV scale in LoNalogy is $M_X=3.4638\times10^{15}\,\mathrm{GeV}$; identifying $M_R=M_X$ introduces no new free parameter.

\emph{Mass unit.}
With $m_D\sim\Lambda_*$ (electroweak scale):
\[
m_\nu^{\rm unit} = \frac{\Lambda_*^2}{M_X} = \frac{(246\,\mathrm{GeV})^2}{3.4638\times10^{15}\,\mathrm{GeV}}
= \frac{6.052\times10^4\,\mathrm{GeV}}{3.4638\times10^{15}\,\mathrm{GeV}}
= 1.747\times10^{-11}\,\mathrm{GeV} = 17.47\,\mathrm{meV}.
\]

\emph{Leading spectrum.}
The $S_3$-covariant Weinberg structure $M_\nu\propto(a-b)I_3+bJ_3$ in the minimal singlet completion ($a=b=1$) gives eigenvalues $(0,0,3m_\nu^{\rm unit})$.
Thus $m_3^{(0)}=3\times17.47=52.4\,\mathrm{meV}$, matching the atmospheric mass-squared difference $\Delta m_{\rm atm}^2\approx(50\,\mathrm{meV})^2$ to within expected $O(1)$ corrections.
The two zero modes are lifted to the solar scale by cusp residues, which is consistent with $\Delta m_{\rm sol}^2\approx(8.6\,\mathrm{meV})^2$.
\end{proof}

\section{Inflation from Cusp Metric}\label{sec:inflation}

\begin{theorem}[Cusp-metric plateau inflation]\label{thm:inflation}
The $x$-field cusp metric $G_{xx}\sim f_T^2/(2\delta^2\log^2(32/\delta))$ produces a double-exponential plateau potential, with normalization-independent slow-roll observables:
\[
\boxed{n_s=1-\frac{2}{N_e}+O\!\left(\frac{1}{N_e^2\log^2N_e}\right),\qquad
r=O\!\left(\frac{f_T^2/M_{\rm Pl}^2}{N_e^2\log^2N_e}\right).}
\]
For $N_e=50$--$60$: $n_s=0.960$--$0.967$, $r\lesssim10^{-4}$ (if $f_T\lesssim M_{\rm Pl}$).
Planck 2018+BK18: $n_s=0.965\pm0.004$, $r<0.036$.
\end{theorem}

\begin{proof}
\emph{Step 1: Cusp metric divergence.}
In the cusp branch $x=1-\delta$, $\delta\downarrow0$, the complete elliptic integral diverges logarithmically:
\[
K\!\left(\frac{1+x}{2}\right)=K\!\left(1-\frac{\delta}{2}\right)\sim\frac{1}{2}\log\frac{32}{\delta},
\qquad K\!\left(\frac{1-x}{2}\right)\to\frac{\pi}{2}.
\]
The field-space metric of the $x$-field action $S_x=\int d^4x\sqrt{-g}\left[-\tfrac12 G_{xx}(\partial x)^2-V(x)\right]$ becomes
\[
G_{xx}(x)\sim\frac{f_T^2}{2\delta^2\log^2(32/\delta)}.
\]

\emph{Step 2: Canonical field.}
The canonical field $\varphi$ defined by $d\varphi=\sqrt{G_{xx}}\,dx$ satisfies
\[
\frac{d\varphi}{d\delta}\approx-\sqrt{G_{xx}}\approx-\frac{f_T}{\sqrt{2}\,\delta\log(32/\delta)}.
\]
Integrating: $\varphi\approx\frac{f_T}{\sqrt{2}}\log\log(32/\delta)$.
Inverting: $\delta=32\exp[-\exp(\sqrt{2}\varphi/f_T)]$.

\emph{Step 3: Plateau potential.}
The potential near the cusp is $V_{79}(x)=V_0-V_1\delta+O(\delta^2)$, giving
\[
U(\varphi)=V_0-32V_1\exp\!\left[-\exp\!\left(\frac{\sqrt{2}\varphi}{f_T}\right)\right]+\cdots.
\]
This is a double-exponential plateau: $U\to V_0$ as $\varphi\to\infty$ (slow roll satisfied).

\emph{Step 4: Slow-roll parameters.}
For large $\varphi$ (many $e$-folds from the end of inflation), the field rolls on the plateau with
\[
\epsilon_V=\frac{M_{\rm Pl}^2}{2}\left(\frac{U'}{U}\right)^2
\sim\frac{M_{\rm Pl}^2 f_T^2\exp[-2\exp(\sqrt{2}\varphi/f_T)]}{V_0^2}
=O\!\left(\frac{f_T^2/M_{\rm Pl}^2}{N_e^2\log^2 N_e}\right),
\]
\[
\eta_V=M_{\rm Pl}^2\frac{U''}{U}\approx-\frac{1}{N_e}+O(N_e^{-2}).
\]
The spectral index and tensor-to-scalar ratio follow:
\[
n_s=1-6\epsilon_V+2\eta_V\approx1+2\eta_V\approx1-\frac{2}{N_e},\qquad
r=16\epsilon_V=O\!\left(\frac{f_T^2/M_{\rm Pl}^2}{N_e^2\log^2N_e}\right).
\]
For $N_e=50$: $n_s=0.960$; for $N_e=60$: $n_s=0.967$.
For $f_T=M_{\rm Pl}$: $r\sim1/(N_e^2\log^2N_e)\sim2\times10^{-4}$.
Both are consistent with Planck 2018+BK18 bounds.
\end{proof}

\section{Consolidated Prediction Table}\label{sec:summary2}

\begin{center}
\begin{tabular}{lccc}
\toprule
Quantity & LoNalogy & Observed & Discrepancy\\
\midrule
$\Omega_{\rm DM}h^2$ & $0.1149$ & $0.1200\pm0.0012$ & $4\%$\\
$\Omega_{\rm DM}/\Omega_{\rm DE}$ & $0.369$ & $0.386$ & $4\%$\\
$\Omega_b h^2$ & $0.02116$ & $0.02237\pm0.00015$ & $5\%$\\
$\eta_B$ & $5.80\times10^{-10}$ & $(6.12\pm0.04)\times10^{-10}$ & $5\%$\\
$m_\nu^{\rm unit}$ & $17.47\,\mathrm{meV}$ & atm.\ $\approx50\,\mathrm{meV}$ & in window\\
$n_s$ & $1-2/N_e$ & $0.965\pm0.004$ & $<1\sigma$\\
$r$ & $\lesssim10^{-4}$ & $<0.036$ & consistent\\
\bottomrule
\end{tabular}
\end{center}

All derived from $(N_{\rm fam},\nu_\br,M_X,Q_{\kker},\Lambda_*)$. Remaining free inputs: $(\rho_0,r_c)$ per halo; $f_T/M_{\rm Pl}$ (inflation amplitude); Yukawa $y$; full 4D gravity UV completion.


%% ============================================================
\section*{Conclusion}
\addcontentsline{toc}{section}{Conclusion}
%% ============================================================

%% ============================================================

Starting from the level-2 modular curve $X(2)$ and the para-complex unit
$\jmath^2=+1$, we have derived:

\begin{enumerate}
\item The SM gauge group $\mathrm{SU}(3)\times\mathrm{SU}(2)\times\mathrm{U}(1)$
  from rank-5 Yukawa/anomaly constraints and $3+2$ holonomy.
\item Exactly three generations from the three cusps of $X(2)$,
  with UV $S_3$ family symmetry.
\item $M_X=3.455\times10^{15}\,\mathrm{GeV}$ and $\tau_p=1.49\times10^{38}\,\mathrm{yr}$
  from the split-ratio formula $M_X=\Lambda_*(r/s)^{1/3}Q_{\ker}^{-14/3}$.
\item $\Lambda_\eff=2.94\times10^{-47}\,\mathrm{GeV}^4$ without free parameters,
  17\% above the PDG 2025 value.
\item A dark fermion mass window $338$--$348\,\mathrm{GeV}$ with
  $\sigma_p^{\rm SI}\sim10^{-48}\,\mathrm{cm}^2$, at the LZ/XENONnT frontier.
\item Analytic halo cores (inner slope $\approx-0.17$) from the $\sech^2$
  coupling geometry.
\end{enumerate}

The most decisive near-future tests are: (i) Hyper-Kamiokande proton
decay search, where non-observation at $10^{35}$--$10^{36}\,\mathrm{yr}$
is predicted; (ii) LZ/XENONnT direct detection at $m_\chi\sim340\,\mathrm{GeV}$
and $\sigma\sim10^{-48}\,\mathrm{cm}^2$.

\section*{Acknowledgements}
\addcontentsline{toc}{section}{Acknowledgements}

The author thanks Claude (Anthropic) and GPT (OpenAI) for assistance with
derivation verification, numerical computation, and manuscript preparation.
Numerical results were verified using Python/NumPy/SciPy
(\texttt{exp\_cosmology.py}, \texttt{exp\_split\_running.py}).

\begin{thebibliography}{99}

\bibitem{YMmassGap}
M.~Nakamura,
\textit{Yang--Mills Existence and Mass Gap: A Constructive Proof via
Elliptic Modular Geometry and Transfer-Poincar\'e Descent},
Zenodo (2026), \href{https://doi.org/10.5281/zenodo.19183838}{doi:10.5281/zenodo.19183838}.

\bibitem{LoNalogyv87}
M.~Nakamura,
\textit{LoNalogy v87.1: Elliptic Modular Framework for Quantum Field Theory
and Millennium Problems},
Zenodo (2026), \href{https://doi.org/10.5281/zenodo.19115338}{doi:10.5281/zenodo.19115338}.

\bibitem{PDG2025}
Particle Data Group (S.~Navas et al.),
\textit{Review of Particle Physics},
Phys.\ Rev.\ D \textbf{110}, 030001 (2024).

\bibitem{LZ2024}
LZ Collaboration,
\textit{Dark Matter Search Results from 4.2 Tonne-Years of Exposure
of the LUX-ZEPLIN (LZ) Experiment},
arXiv:2410.17036 (2024).

\bibitem{HyperK}
Hyper-Kamiokande Collaboration,
\textit{Search for proton decay via $p\to e^+\pi^0$ and $p\to\mu^+\pi^0$},
arXiv:2010.16098 (2020).

\bibitem{DESI2025}
DESI Collaboration,
\textit{DESI 2025 Results},
in preparation.

\end{thebibliography}

\end{document}
